% Tutorial 5: USDC Depeg - ECOM215
\documentclass[11pt]{article}
\usepackage{fullpage}
\usepackage[left=1in,top=1in,right=1in,bottom=1in,headheight=3ex,headsep=3ex]{geometry}
\usepackage{graphicx}
\usepackage{float}
\usepackage{booktabs}
\usepackage{enumitem}

\usepackage[sc]{mathpazo}
\linespread{1.05}
\usepackage[T1]{fontenc}

\usepackage{fancyhdr,lastpage}
\pagestyle{fancy}
\usepackage{hyperref}
\usepackage{lastpage}

\usepackage{array, xcolor}
\usepackage{colortbl}
\definecolor{clemsonorange}{HTML}{EA6A20}
\definecolor{ImperialBlue}{RGB}{0, 62, 116}
\definecolor{AccentGreen}{RGB}{46, 139, 87}
\hypersetup{colorlinks,breaklinks,linkcolor=clemsonorange,urlcolor=clemsonorange,anchorcolor=clemsonorange,citecolor=black}

\lhead{}
\chead{}
\rhead{\footnotesize ECOM215: Tutorial 5 --- USDC Depeg}
\lfoot{}
\cfoot{\small \thepage/\pageref*{LastPage}}
\rfoot{}

\title{Tutorial 5: The USDC Depeg\\[0.3em] \Large When ``Safe'' Isn't Safe\\[1em] \normalsize ECOM215: Blockchain Economics and Digital Assets}
\author{Week 6 $\vert$ Stablecoins and Central Bank Digital Currencies}
\date{Semester B, 2025/2026}

\begin{document}

\maketitle

\vspace{1em}
\begin{center}
\fbox{\parbox{0.9\textwidth}{\centering\textbf{CASE BRIEF FOR STUDENTS}\\[0.3em] Please read before the tutorial. Estimated reading time: 15 minutes.}}
\end{center}

\section*{Background}

By early 2023, USD Coin (USDC) had established itself as the ``safe'' stablecoin. Unlike Tether (USDT), which faced persistent questions about its reserves, USDC was issued by Circle---a US-regulated company with monthly attestations from major accounting firms. Its reserves were held in cash and short-dated US Treasury bills at regulated financial institutions.

USDC was the backbone of DeFi. It served as the primary stablecoin for lending protocols, DEX liquidity pools, and cross-border payments. Its market capitalisation exceeded \$40 billion. It was, by design, boring and reliable.

Then, on a Friday evening in March 2023, USDC dropped to \$0.87.

\section*{The Silicon Valley Bank Collapse}

\textbf{Silicon Valley Bank (SVB)} was a mid-sized US bank specialising in technology startups and venture capital. It was the 16th largest bank in the United States.

\medskip
\noindent\textbf{What went wrong at SVB:}
\begin{itemize}[leftmargin=*]
\item SVB held large amounts of long-dated government bonds and mortgage-backed securities
\item When interest rates rose sharply in 2022, these bonds lost value (bond prices fall when rates rise)
\item SVB had \$17 billion in unrealised losses on its bond portfolio
\item On March 8, 2023, SVB announced it needed to raise \$2.25 billion to shore up its balance sheet
\item Depositors panicked; \$42 billion was withdrawn in a single day
\item On March 10 (Friday), regulators seized SVB---the second-largest bank failure in US history
\end{itemize}

\section*{The USDC Connection}

On the evening of Friday, March 10, Circle disclosed that \textbf{\$3.3 billion of USDC's cash reserves}---approximately 8\% of total reserves---were held at Silicon Valley Bank.


\medskip
\noindent\textbf{The immediate market reaction:}
\begin{itemize}[leftmargin=*]
\item USDC began trading below \$1 within minutes of the announcement
\item By Saturday morning, USDC had dropped to \$0.87--\$0.90 on major exchanges
\item Crypto markets broadly sold off; Bitcoin dropped 8\%
\item DeFi protocols faced stress as USDC collateral lost value
\item DAI, which uses USDC as partial backing, also briefly lost its peg
\end{itemize}

\medskip
\noindent\textbf{The key question}: If Circle couldn't access \$3.3 billion of its reserves, could it honour redemptions at \$1?

\section*{48 Hours of Uncertainty}

\noindent\textbf{Friday evening, March 10:}
\begin{itemize}[leftmargin=*]
\item Circle's disclosure was transparent but alarming
\item Traditional banking was closed for the weekend---no clarity on whether SVB deposits would be recovered
\item USDC redemptions require banking infrastructure; weekend = no redemptions possible
\item Holders could only sell on secondary markets, driving prices down
\end{itemize}

\medskip
\noindent\textbf{Saturday, March 11:}
\begin{itemize}[leftmargin=*]
\item USDC traded between \$0.87 and \$0.95 across different venues
\item Arbitrage was impossible without redemption access
\item Rumours and speculation dominated crypto social media
\item Some DeFi protocols paused operations; others faced liquidation pressure
\end{itemize}

\medskip
\noindent\textbf{Sunday, March 12:}
\begin{itemize}[leftmargin=*]
\item US Treasury, Federal Reserve, and FDIC issued joint statement
\item \textbf{All SVB depositors would be made whole}---including uninsured deposits above \$250,000
\item Circle confirmed it would have full access to its SVB funds
\item USDC began recovering toward \$1
\end{itemize}

\medskip
\noindent\textbf{Monday, March 13:}
\begin{itemize}[leftmargin=*]
\item Banking resumed; Circle processed redemptions normally
\item USDC returned to \$1 peg
\item Crisis over---but lessons remained
\end{itemize}

\section*{Why Did the ``Safe'' Stablecoin Depeg?}

\textbf{1. Counterparty risk from banking partners}
\begin{itemize}[leftmargin=*]
\item USDC's reserves were held in banks, not directly in Treasuries
\item Bank deposits above \$250,000 are not FDIC-insured
\item Circle had \$3.3B at SVB---all uninsured
\end{itemize}

\medskip
\noindent\textbf{2. Weekend liquidity mismatch}
\begin{itemize}[leftmargin=*]
\item Crypto markets trade 24/7; banks operate on business days
\item Redemptions were impossible over the weekend
\item Arbitrage (buy cheap USDC, redeem for \$1) couldn't function
\item Market price reflected uncertainty, not fundamentals
\end{itemize}

\medskip
\noindent\textbf{3. Transparency as a double-edged sword}
\begin{itemize}[leftmargin=*]
\item Circle disclosed the SVB exposure promptly---arguably the right thing to do
\item But the disclosure triggered the panic
\item Tether (less transparent) saw inflows as a ``safe haven'' during the crisis
\end{itemize}

\section*{The Government Backstop}

The US government's decision to guarantee all SVB deposits---not just the FDIC-insured \$250,000---was extraordinary.

\medskip
\noindent\textbf{The rationale:}
\begin{itemize}[leftmargin=*]
\item Prevent contagion to other regional banks
\item Protect payroll and operations of thousands of companies
\item Maintain financial stability
\end{itemize}

\medskip
\noindent\textbf{The controversy:}
\begin{itemize}[leftmargin=*]
\item Was this a bailout? (Officials insisted it was not---shareholders and bondholders lost everything)
\item Moral hazard: Does this encourage risky behaviour by banks and depositors?
\item Should crypto companies benefit from government safety nets they often criticise?
\end{itemize}

\medskip
\noindent\textbf{For USDC}: The government backstop saved Circle from a potential crisis. Without it, USDC might have faced a genuine solvency problem.

\section*{Aftermath and Changes}

\textbf{Circle's response:}
\begin{itemize}[leftmargin=*]
\item Diversified banking relationships (no single bank holds large share)
\item Increased allocation to Treasury bills (held at BNY Mellon)
\item Reduced reliance on cash deposits
\item Launched ``Circle Reserve Fund'' with BlackRock
\end{itemize}

\medskip
\noindent\textbf{Market impact:}
\begin{itemize}[leftmargin=*]
\item USDC market cap declined from \$43B to \$25B over following months
\item Tether (USDT) market cap increased---perceived as less connected to US banking
\item Some DeFi protocols diversified stablecoin holdings
\end{itemize}

\medskip
\noindent\textbf{Regulatory implications:}
\begin{itemize}[leftmargin=*]
\item Renewed focus on stablecoin reserve requirements
\item Questions about whether stablecoins should have direct central bank accounts
\item Debate over 24/7 redemption requirements
\end{itemize}

\section*{Key Facts Summary}

\begin{center}
\begin{tabular}{ll}
\toprule
\textbf{Item} & \textbf{Detail} \\
\midrule
Date & March 10--13, 2023 \\
USDC market cap & $\sim$\$43 billion \\
SVB exposure & \$3.3 billion ($\sim$8\% of reserves) \\
USDC low price & $\sim$\$0.87 \\
Duration of depeg & $\sim$48 hours \\
Resolution & US government guaranteed all SVB deposits \\
\bottomrule
\end{tabular}
\end{center}

\section*{Questions to Consider}

\begin{enumerate}[leftmargin=*]
\item Was Circle's immediate disclosure the right decision? What are the trade-offs between transparency and stability?

\item Should stablecoin issuers be allowed to hold reserves in commercial banks? What alternatives exist?

\item The government backstop saved USDC. Is it appropriate for crypto to benefit from traditional financial safety nets?

\item How should stablecoin reserves be structured to avoid this type of risk?

\item Does this case change your view of ``fiat-backed'' stablecoins as the ``safe'' option?
\end{enumerate}

\section*{Further Reading (Optional)}

\begin{itemize}[leftmargin=*]
\item Circle's March 10 and March 11 public statements
\item Joint statement by Treasury, Federal Reserve, and FDIC (March 12, 2023)
\item ``Lessons from the USDC Depeg''---Chainalysis blog
\end{itemize}

\newpage

%==============================================================================
% PART B: TA DISCUSSION GUIDE
%==============================================================================

\section*{Session Timeline}

\begin{center}
\begin{tabular}{ll}
\toprule
\textbf{Time} & \textbf{Activity} \\
\midrule
0:00--0:08 & Context setting: summarise the SVB collapse and USDC connection \\
0:08--0:20 & Discussion Question 1: Transparency trade-offs \\
0:20--0:32 & Discussion Question 2: Reserve structure \\
0:32--0:44 & Discussion Question 3: Government backstops \\
0:44--0:52 & Discussion Question 4: What is ``safe''? \\
0:52--1:00 & Synthesis and key takeaways \\
\bottomrule
\end{tabular}
\end{center}

\section*{Discussion Questions with Guidance}

\subsection*{Question 1: Was immediate disclosure the right call?}

\textit{``Circle disclosed its SVB exposure on Friday evening, triggering the depeg. Was this the right decision? Would delaying until Monday have been better?''}

\medskip
\noindent\underline{Arguments for immediate disclosure:}
\begin{itemize}[leftmargin=*]
\item Legal and ethical obligation to inform holders of material risk
\item If news leaked through other channels, trust would be destroyed
\item Transparency is Circle's brand differentiator---silence would undermine it
\item Markets will eventually learn; better to control the narrative
\end{itemize}

\medskip
\noindent\underline{Arguments for delaying:}
\begin{itemize}[leftmargin=*]
\item Weekend timing meant no redemptions possible---panic with no relief valve
\item By Monday, government response might have been known, reducing panic
\item Disclosure without solution can cause unnecessary harm
\item Ironically, Tether's opacity shielded it from scrutiny
\end{itemize}

\medskip
\noindent\textbf{Key insight}: Transparency is valuable but creates its own dynamics. In a 24/7 market with weekend banking closures, timing of disclosure matters enormously. There's no clearly ``correct'' answer here.

\subsection*{Question 2: How should stablecoin reserves be structured?}

\textit{``USDC was considered safe because it held cash and Treasuries. But the cash was in a bank that failed. How should reserves be structured to avoid this?''}

\medskip
\noindent\underline{Options to discuss:}
\begin{itemize}[leftmargin=*]
\item \textbf{100\% Treasury bills}: Safe from bank failure, but requires custody infrastructure
\item \textbf{Direct central bank accounts}: Safest possible, but Fed doesn't offer this to non-banks
\item \textbf{Multiple bank diversification}: Spreads risk, but increases complexity
\item \textbf{Government money market funds}: Low risk, regulated, liquid
\item \textbf{Full reserve banking license}: Become a bank, but heavy regulation
\end{itemize}

\medskip
\noindent\underline{Trade-offs:}
\begin{itemize}[leftmargin=*]
\item Safety vs.\ yield (Treasuries pay interest; spreading across banks is operationally complex)
\item Centralisation vs.\ resilience (fewer partners = simpler but riskier)
\item Regulatory access (some safe options require licenses stablecoin issuers don't have)
\end{itemize}

\medskip
\noindent\textbf{Key insight}: There is no perfectly safe option within existing infrastructure. Even ``risk-free'' assets require custody, which introduces counterparty risk. This is why some argue stablecoins need direct central bank access or should become regulated banks.

\subsection*{Question 3: Should crypto benefit from government safety nets?}

\textit{``The US government's decision to guarantee all SVB deposits saved USDC. Is it appropriate for crypto---which often positions itself as an alternative to traditional finance---to benefit from government backstops?''}

\medskip
\noindent\underline{Arguments that it's appropriate:}
\begin{itemize}[leftmargin=*]
\item Circle followed the rules; it used the banking system as designed
\item Backstop was for all depositors, not specifically for crypto
\item Stablecoins serve real economic functions (payments, DeFi); their failure would harm people
\item Crypto companies pay taxes and employ people---they're part of the economy
\end{itemize}

\medskip
\noindent\underline{Arguments that it's problematic:}
\begin{itemize}[leftmargin=*]
\item Crypto's value proposition includes avoiding government intervention
\item Moral hazard: Crypto companies may take more risk expecting bailouts
\item If crypto wants to be ``outside the system,'' it shouldn't benefit from system protections
\item Creates unfair advantage over truly decentralised alternatives
\end{itemize}

\medskip
\noindent\textbf{Key insight}: This exposes a fundamental tension. Stablecoins like USDC are deeply integrated with traditional finance---they're not really ``outside the system.'' The question is whether this integration should be formalised through regulation or remains an awkward hybrid.

\subsection*{Question 4: What does ``safe'' mean for stablecoins?}

\textit{``Before March 2023, USDC was considered the safest major stablecoin. Does this event change your assessment? What should investors look for?''}

\medskip
\noindent\underline{Points to explore:}
\begin{itemize}[leftmargin=*]
\item ``Safe'' is relative, not absolute---all financial instruments have risks
\item USDC's risk was counterparty (bank failure), not fraud or design flaw
\item The depeg was brief and fully resolved---arguably the system worked
\item But: Resolution depended on extraordinary government intervention
\item Transparency helped Circle recover trust, even if it caused short-term pain
\end{itemize}

\medskip
\noindent\underline{What to look for in stablecoin safety:}
\begin{itemize}[leftmargin=*]
\item Reserve composition and quality
\item Custodian and banking partner diversification
\item Regulatory status and jurisdiction
\item Track record through stress events
\item Transparency and audit quality
\item Redemption mechanisms and liquidity
\end{itemize}

\medskip
\noindent\textbf{Key insight}: No stablecoin is perfectly safe. The USDC case shows that even well-designed, transparent, regulated stablecoins face risks from the financial infrastructure they depend on. Due diligence requires understanding the full chain of dependencies.

\subsection*{Extension Question (if time permits)}

\textit{``Tether gained market share during and after the USDC crisis, despite being less transparent. What does this tell us about market preferences?''}

\medskip
This raises interesting questions about the value of transparency. USDT's opacity meant it didn't have to disclose exposure; investors fled to it as a ``known unknown'' rather than USDC's ``known risk.'' Is this rational? What does it imply for optimal disclosure policy?

\vfill
\begin{flushright}
\textit{End of Tutorial 5 Materials}
\end{flushright}

\end{document}